\subsection{Integer divisions and actual trace operators}
\begin{proposition}[Arithmetic realization of the coefficient rules]\label{ex:app-coefficient-rules}
Use the rational full-two-torsion graph conventions of
Definitions~\ref{ex:sec7-typed-graph} and~\ref{ex:sec7-symbol}.
For an even source of Proposition~\ref{ex:app-even-phase}, use its
fixed minimum depth. For an odd source, first fix a finite operator
budget $|U|\le W$, take $R\ge2$, and choose an actual stage satisfying
the depth inequalities of Proposition~\ref{ex:app-odd-minimum}.
Then the integer divisions (O.9)--(O.11) are valid at every Fourier
index, with the stated truncation for the odd source.
The fresh coefficient-one reading is
$\overline{(t(\operatorname{Fr}_q)\theta H_U)[1]}=\mathcal A(Uq)$.
On all tested nonempty configurations, these actual symbols depend
only on their typed graph and satisfy the rules of
Lemma~\ref{ex:sec7-coefficient-rules}: a pair difference
contracts the endpoints and adds their types and every external
incidence; whole-prime omission uses the same source, phase, depth
and stage; and an isolated primary vertex of neutral type and
neutral external incidences makes the symbol zero.
Every contraction has nonempty output. Stationary character tests
are carried along in the same product rule. This statement asserts
neither values at untested odd-stage weights nor an identity at an
empty contraction output.
\end{proposition}
\begin{proof}\leavevmode
All new primes below avoid the full coefficient support $S$,
including every late sieve and the current $\ell_0$.
Here $S$ is the actual coefficient support.  The coefficients of
$g$ vanish at all indices not prime to $S$.
For the odd source this is forced at $2,N,h_*,\ell_0$ by the
local Schwartz support and the extraction at $|h_*|$.
At every additional coefficient-sieve prime it is imposed by an
exact zero-or-retain Fourier sieve, which commutes with the good
recurrence outside the enlarged support.  This is the vanishing
used in the condition $(m,S)=1$ in the unary series below.
Later stationary prime-selection tests may enrich the common
label alphabet without resieving the intrinsic odd source;
such tests do not enlarge this coefficient-support condition.
Use the minimum-depth inequalities proved in Appendix~\ref{app:endpoints}.
The general bookkeeping assigns weight one to a split prime of
$\Q(E[2])$ and one-half to a prime with transposition Frobenius, called
simple. In our actual full-two-torsion family every prime is split,
$w(U)=|U|$, and the depth $\delta$ is an integer. Thus only matching
parity, meaning $e_U=0$, is needed here; the displayed general formula
also records where the unused simple-prime term would enter.
For a finite set $U$, put
\[
 n_U=\lceil\delta+w(U)\rceil,\quad e_U=n_U-\delta-w(U),\qquad
 B_U=2^{-n_U}\prod_{p\in U}(U_p-\eta(p)V_p)g,
\]
where $U_p(\sum a(n)q^n)=\sum a(pn)q^n$ and $V_pg(\tau)=g(p\tau)$.
These are operations on the actual form $g$. We verify the level
at a new odd prime using the classical theta multiplier. Write the
old level as $4A$, with $p\nmid4A$, and put
$\beta=\left(\begin{smallmatrix}p&0\\0&1\end{smallmatrix}\right)$.
For $\gamma=\left(\begin{smallmatrix}a&b\\c&d\end{smallmatrix}\right)
\in\Gamma_0(4Ap)$ one has
\[
 \beta\gamma\beta^{-1}=
       \begin{pmatrix}a&pb\\c/p&d\end{pmatrix}\in\Gamma_0(4A).
\]
The square-root factor in the theta multiplier is unchanged on
evaluating the latter matrix at $p\tau$. The quotient of its
arithmetic factor by that of $\gamma$ is
$(c/p\,/d)/(c/d)=(p/d)$. Cubing the multiplier leaves this
quadratic factor unchanged. Hence $V_pg$ has level $4Ap$ and
character $\psi(p/\cdot)$.

At this level the normalized sum
\[
 U_{p^2}h(\tau)=\frac1{p^2}\sum_{r\bmod p^2}
                             h((\tau+r)/p^2)
\]
preserves the same weight and character. Indeed, put
$\alpha_r=\left(\begin{smallmatrix}1&r\\0&p^2\end{smallmatrix}\right)$.
For $\gamma\in\Gamma_0(4Ap)$ choose
$r'\equiv(b+rd)/(a+rc)\pmod{p^2}$. Then
\[
 \alpha_r\gamma=\gamma'_r\alpha_{r'},\qquad
 \gamma'_r=
 \begin{pmatrix}
  a+rc&(b+rd-(a+rc)r')/p^2\\
  p^2c&d-cr'
 \end{pmatrix}\in\Gamma_0(4Ap).
\]
The map $r\mapsto r'$ is a permutation, since $a+rc$ is a unit
modulo $p$ and the fractional linear map has determinant one.
Writing $d'=d-cr'$, its nebentypus factor is unchanged since
$d'\equiv d\pmod{4Ap}$. Its theta arithmetic factor is unchanged
as well: $d'\equiv d\pmod c$, $d'\equiv d\pmod4$, and
$(p^2c/d')=(c/d')=(c/d)$. These are the usual Kronecker-symbol
identities for the theta multiplier on $\Gamma_0(4)$.
Its square-root factor, evaluated at $(\tau+r')/p^2$, is
$(c\tau+d)^{1/2}$. Thus the sum transforms with the same
theta multiplier cubed and the same character. The formulas also
hold for the triangular matrices, and holomorphy at the cusps is
preserved by the finite rational translates.
On Fourier expansions $U_{p^2}V_p=U_p$. Consequently both terms
of $U_p-\eta(p)V_p$ have level exponent at most one at $p$ and
character $\psi(p/\cdot)$. Iteration introduces no higher exponent
at previously added primes.

The same computation applies to the unary corrections. Start with
the odd-character unary theta series at the old support, including
its fixed coprimality sieve. Its $J$-series is its transform by
$V_J$, for squarefree $J\subset U$. A prime $p\mid J$ therefore
has exponent at most one and the extra character $(p/\cdot)$;
a prime $p\nmid J$ is absent from its original level. Multiplication
by the ordinary theta series introduces no odd level prime. This
proves the precise level and oldform alternatives used in (O.12).


